
\section{正则性与域上光滑}

\trRegularity{}定义中的关键想法，已包含在定理~12.3.12 之中（\trZariski{}\trTangentSpace{}的维数至少是\trLocalRing{}的维数）。

\textbf{13.2.1.}\ \emph{\trhDefinition{}。}\quad 若定理~12.3.12 中取等号，则称 $A$ 是一个\textbf{\trRegular\trLocalRing}（regular local ring）。特别地，\emph{\trRegular\trLocalRing{}按定义是\trNoetherian{}的}。（若一个\trNoetherian{}环 $A$ 在它的所有\trPrimeIdeal{}处都正则——也就是说，若对 $A$ 的所有\trPrimeIdeal{} $\mathfrak{p}$，$A_{\mathfrak{p}}$ 都是\trRegular\trLocalRing——则称 $A$ 是一个\textbf{\trRegular{}环}（regular ring）。这个术语我们基本不去使用。事实~13.8.2 令人安心地蕴涵：\trRegular\trLocalRing{}的确是\trRegular{}环。）一个\trLocally\trNoetherian{}概形 $X$ 称为在一点 $p$ 处是\textbf{正则}的（regular），若\trLocalRing{} $\OO_{X,p}$ 是正则的。（「\trNonsingular{}的」（nonsingular）一词也常被使用，特别是当 $X$ 是域上的\trFiniteType{}概形时；但为了口径一致，我们全篇使用「正则」。）否则称它在这一点处是\textbf{\trSingular{}的}（singular），并称这个点是一个\textbf{\trSingularPoint}（singularity）。（「非正则的」（nonregular）一词也在使用，但同样为了口径一致，我们将使用「\trSingular{}的」——尽管这个选择，与我们宁取「正则」而不取「\trNonsingular」的选择是\emph{不相一致}的。这种不相一致的一个可能的动机是：「正则局部环」与「\trSingularPoint」都是标准术语，所以总有个时刻我们不得不做出选择。）一个概形是\textbf{正则}的（regular）（或\textbf{\trNonsingular{}的}（nonsingular）），若它在所有点处都正则；否则它是\textbf{\trSingular{}的}（singular）（也就是说，它在\emph{至少一个}点处\trSingular——若它有一个\trSingularPoint）。

\textbf{13.2.A.（练习）}\quad 证明：维数为 $0$ 的\trNoetherian\trLocalRing{}正则，当且仅当它是一个域。

希望你将渐渐信服：这就是概形的「\trSmoothness」的正确观念。值得注意的是，\trKrull{}出于纯代数的理由引入了\trRegular\trLocalRing{}的观念，那比\trZariski{}于 1947 年认识到它是几何中的一个基本观念，还要早上一些时候。

\textbf{13.2.B.（练习）}\quad 设 $(A,\mathfrak{m})$ 是维数为 $n>0$ 的\trRegular\trLocalRing，而 $f\in A$。证明：$A/(f)$ 是维数为 $n-1$ 的\trRegular\trLocalRing，当且仅当 $f\in\mathfrak{m}\setminus\mathfrak{m}^2$。（提示：切空间版的\trKrull{}\trPrincipalIdealTheorem，\mbox{练习~13.1.B}。）

\textbf{13.2.C.（练习：\trRegularity{}的切片判别法）}\quad 设 $X$ 是一个\trNoetherian{}概形（例如一个\trVariety），$D$ 是 $X$ 上的一张\trEffectiveCartierDivisor（定义~9.5.1），而 $p\in D$。证明：若 $p$ 是 $D$ 的正则点，则 $p$ 是 $X$ 的正则点。（提示：切空间版的\trKrull{}\trPrincipalIdealTheorem，\mbox{练习~13.1.B}。）

\textbf{13.2.D.（容易的练习：操练概念）}\quad 设 $(A,\mathfrak{m})$ 是维数为 $d$ 的\trRegular\trLocalRing，而 $f_1,\dots,f_d\in\mathfrak{m}$，它们在\trVectorSpace{} $\mathfrak{m}/\mathfrak{m}^2$ 中的像构成一组基。证明：它们是 $(A,\mathfrak{m})$ 的一个\trSystemOfParameters（定义~12.3.11）。（它们是「特别好」的参数。）

\textbf{13.2.2.\ \trRegularity{}的\trJacobianCriterion。}

\trFiniteType{}的 $k$-\trScheme{}局部地都形如 $\Spec k[x_1,\dots,x_n]/(f_1,\dots,f_r)$。当 $k=\overline{k}$ 时，\trRegularity{}的\trJacobianCriterion（\mbox{练习~13.2.E}）给出一个可以动手操作的方法，利用方程 $f_1,\dots,f_r$ 来检验\trClosedPoint{}处的\trSingular{}性。

\textbf{13.2.E.（重要练习：\trRegularity{}的\trJacobianCriterion——有了\mbox{练习~13.1.I}便容易）}\quad 设 $X=\Spec k[x_1,\dots,x_n]/(f_1,\dots,f_r)$ 具有\trPureDimension{} $d$。证明：$k$-\trValuedPoint{} $p\in X$ 是正则的，当且仅当雅可比\trCorank{}函数在 $p$ 处取值为 $d$。

\textbf{13.2.F.（容易的练习）}\quad 设 $k=\overline{k}$。证明：$\mathbb{A}^n_k$ 中超曲面 $f(x_1,\dots,x_n)=0$ 的\trSingular\trClosedPoint{}由方程组
\[
f=\frac{\partial f}{\partial x_1}=\cdots=\frac{\partial f}{\partial x_n}=0
\]
给出。（转述：$f=0$ 的\trSingular{}点，就是 $f$ 的梯度消失之处。这一点也不令人震惊。）

\textbf{13.2.3.\ 域上的\trSmoothness。}

\trJacobianCriterion{}似乎有两个严重的缺陷。对 $\overline{k}$ 上的\trFiniteType{}概形，这条判别法给出了\trRegularity{}的一个必要条件，但其是否充分并不显然——我们还需要在非\trClosedPoint{}处也检验\trRegularity。我们可以通过费力证明事实~13.8.2 来证得充分性，该事实蕴涵：非\trClosedPoint{}也必是正则的。第二个缺陷是：这条判别法要求 $k$ 是\trAlgClosedField。这些问题让人以为：使用导数与雅可比矩阵的老式想法，配不上花哨的现代\trRegularity{}观念。但这样想就错了——毛病出在\trRegularity{}这个概念身上。有一个更好的观念：域上的\trSmoothness。更妙的是，这个观念还能推广到概形的\trSmooth\trMorphism{}观念（将在 \S13.6 中讨论，并将在第~24.8 章中再次讨论），它在一切可能的方面都表现良好（包括 \S8.1 的意义上）。这是又一个迹象，表明我们原以为属于对象的一些性质（「绝对观念」），其实应当被理解为态射的性质（「相对观念」）。我们已经懂得足够多，足以（不完美地、然而正确地）定义一个概形何时是\textbf{$k$-光滑}的，或\textbf{在 $k$ 上光滑}的了。

\textbf{13.2.4.}\ \emph{\trhDefinition{}。}\quad 一个 $k$-\trScheme{}称为 \textbf{维数 $d$ 的 $k$-光滑}（$k$-smooth of dimension $d$）概形，或 \textbf{在 $k$ 上维数为 $d$ 的光滑}（smooth of dimension $d$ over $k$）概形，若它具有\trPureDimension{} $d$，并且存在一个由\trAffine\trOpenSubset{} $\Spec k[x_1,\dots,x_n]/(f_1,\dots,f_r)$ 组成的\trCover，使\trJacobianMatrix{}在\emph{所有}点（而不仅在\emph{\trClosed{}点}）处的\trCorank{}为 $d$。（特别地，光滑概形不言而喻地是\trLocally\trFiniteType{}的。）一个 $k$-\trScheme{}称为\textbf{在 $k$ 上光滑}的（smooth over $k$），若它在某个维数下光滑。上下文清楚时，$k$ 常被省略。由 \S13.1.7 中发展的雅可比\trCorank{}函数理论，对任何\trPureDimension{}为 $d$ 的\trVariety，这一点可以在\emph{任何一个}\trAffine\trOpenCover{}上加以检验。

\textbf{13.2.G.（练习：最初的例子）}
\begin{enumerate}[label=(\alph*)]
\item 证明：对任何 $n$ 与任何域 $k$，$\mathbb{A}^n_k$ 都是光滑的。对哪些特征，$\mathbb{P}^2_k$ 中的曲线 $y^2z=x^3-xz^2$ 是光滑的（对照\mbox{练习~13.3.C}）？
\item 设 $f\in k[x_1,\dots,x_n]$ 是一个多项式，使得方程组
\[
f=\frac{\partial f}{\partial x_1}=\cdots=\frac{\partial f}{\partial x_n}=0
\]
在 $\overline{k}$ 中无解。证明：$\mathbb{A}^n_k$ 中的超曲面 $f=0$ 是光滑的。（把这道题与\mbox{练习~13.2.F}相比较，后者多出 $k=\overline{k}$ 这条假设。）
\end{enumerate}

\textbf{13.2.H.（练习：簇的光滑性「与基域的扩张无关」）}\quad 设 $X$ 是一个\trLocally\trFiniteType{}的 $k$-\trScheme，而 $k\subset\ell$ 是一个\trFieldExtension。证明：$X$ 在 $k$ 上光滑，当且仅当 $X\times_{\Spec k}\Spec\ell$ 在 $\ell$ 上光滑。

下一道练习表明：我们只需检验\trClosedPoint{}，从而与经典几何搭上了关系。

\textbf{13.2.I.（练习）}\quad 证明：若 $X=\Spec k[x_1,\dots,x_n]/(f_1,\dots,f_r)$ 的\trJacobianMatrix{}在所有\emph{\trClosed{}点}处的\trCorank{}为 $d$，则它在\emph{所有}点处的\trCorank{}为 $d$。（提示：\trJacobianMatrix{}\trCorank{}为 $d$ 的那个\trLocus，可以用某些显式矩阵的行列式的消失与非消失来描述。）

\textbf{13.2.5.}\ \emph{为\trRegularity{}辩护。}\quad 我们方才替\trSmoothness{}慷慨陈词，但也应当说清楚：\trRegularity{}仍然非常有用。例如，它是唯一在\trMixedCharacteristic（mixed characteristic）下说得通的概念。特别地，$\mathbb{Z}$ 在它的各个点处都正则（它是一个\trRegular{}环），而更一般地，\trDVR{}是\trRegular{}环的极其有用的例子。

\textbf{13.2.6.\ \trRegularity{}对\trSmoothness。}

\textbf{13.2.J.（练习）}\quad 设 $X$ 是\trAlgClosedField{} $k=\overline{k}$ 上\trPureDimension{}为 $d$ 的一个\trLocally\trFiniteType{}概形。证明：$X$ 在它的诸\trClosedPoint{}处正则，当且仅当它是光滑的。（我们很快就会知道：对\trLocally\trFiniteType{}的 $\overline{k}$-概形，闭点处的正则性与处处正则是一回事，见定理~13.8.3。）证明「正则蕴涵光滑」的提示：用\trJacobianCriterion{}证明\trJacobianMatrix{}的\trCorank{}在 $X$ 的诸\trClosedPoint{}处为 $d$，然后用\mbox{练习~13.2.I}。

更一般地，若 $k$ 是完全域（例如 $\operatorname{char}k=0$ 或 $k$ 是\trFinite{}域），则\trSmoothness{}与\trClosedPoint{}处的\trRegularity{}是一回事（见\mbox{练习~13.2.R}）。再更一般地，我们日后将证明下述事实。现在就把它摆出来，是因为远在最终证明它之前，它就能让一大批陈述变得干净利落。（不会有循环论证。）

\textbf{13.2.7.\ 光滑性--正则性比较定理。}\ ——
\begin{enumerate}[label=(\alph*)]
\item \emph{若 $k$ 是完全域，则每个正则的\trFiniteType{} $k$-\trScheme{}都在 $k$ 上光滑。}
\item \emph{每个光滑的 $k$-\trScheme{}都是正则的（对完全性不作任何假设）。}
\end{enumerate}

(a) 将在\mbox{练习~13.2.R}中证明。(b) 将在 \S24.8.9 中证明。定理~13.2.7(b) 要到那么遥远的将来才获证明，并不意味着我们真的得等那么久。我们是有意把证明推迟到手上有了能替我们做掉大部分工作的工具之时。如果你希望立刻得到一点洞见，这里是 (b) 的论证梗概；读完本章之后，你甚至可以试着把它实现出来。我们很快将证明（\mbox{练习~13.3.P}）：对一切域 $k$，$\mathbb{A}^n_k$ 都是正则的。一旦知道了\trEtale{}\trMorphism{}的定义，我们就会认识到：每个光滑\trVariety{}局部地都容有一个到 $\mathbb{A}^n_k$ 的\trEtale{}\trMorphism（观察~13.6.4）。那时你便能证明：若 $\pi\colon X\to Y$ 是\trLocally\trNoetherian{}概形之间的\trEtale{}\trMorphism，而 $p\in X$，则只要 $\pi(p)$ 正则，$p$ 就正则。这将在\mbox{练习~24.8.E}中获得证明，不过你在读完平展性的定义之后就去证明它，也毫不失当。

\textbf{13.2.8.}\ \emph{\trhCaveat{}：\trRegularity{}不蕴涵\trSmoothness。}\quad 若 $k$ 不是完全域，则\trRegularity{}\emph{不}蕴涵\trSmoothness，下面的例子即为明证。令 $k=\mathbb{F}_p(u)$，并考虑超曲面 $X=\Spec k[x]/(x^p-u)$。此时 $k[x]/(x^p-u)$ 是一个域，故而是正则的。但若 $f(x)=x^p-u$，则 $f(u^{1/p})=\frac{df}{dx}(u^{1/p})=0$，于是\trJacobianCriterion{}失效——$X$ \emph{在 $k$ 上}并不光滑。（千万别忘记：\trSmoothness{}需要一个域的选取——它是一个「相对」观念；日后我们将在 \S13.6 中定义任意概形上的\trSmoothness。）

\textbf{13.2.9.}\quad 倘若嫌上一个例子太「小」而不足以启人深思（因为所论概形在\emph{另一个}域上\emph{确}实光滑，即在 $k[x]/(x^p-u)$ 上光滑），这里再给出一个。同前取 $k=\mathbb{F}_p(u)$，但设 $p>2$，并考虑曲线 $\Spec k[x,y]/(y^2-x^p+u)$。那么\trClosedPoint{} $(y,x^p-u)$ 是正则的，却不光滑。

因此，你不要把「正则」与「光滑」当作同义词混用。

\textbf{13.2.10.\ \trRegular\trLocalRing{}都是\trIntegralDomain。}

凭几何直觉你或许会预期：概形在一个「光滑」的点处是「\trLocally\trIrreducible」的。用代数语言来说：

\textbf{13.2.11.\ \trhTheorem{}。}\ ——\ \emph{设 $(A,\mathfrak{m},k)$ 是维数为 $n$ 的\trRegular\trLocalRing。那么 $A$ 是一个\trIntegralDomain。}

在证明它之前（\S13.2.14），我们先给出几条推论。

\textbf{13.2.K.（练习）}\quad 设 $p$ 是一个\trLocally\trNoetherian{}概形 $X$ 的一个正则点。证明：$X$ 只有一个\trIrreducibleComponent{}经过 $p$。

\textbf{13.2.L.（容易的练习）}\quad 证明：非空的\trRegular\trNoetherian{}概形是\trIrreducible{}的，当且仅当它是\trConnected{}的。

\textbf{13.2.M.（重要练习：\trRegular{}概形之中的\trRegular{}概形是\trRegularEmbedding）}
\begin{enumerate}[label=(\alph*)]
\item 设 $(A,\mathfrak{m},k)$ 是维数为 $n$ 的\trRegular\trLocalRing，而 $I\subset A$ 是一个理想，它切出一个维数为 $d$ 的\trRegular\trLocalRing。令 $r=n-d$。证明：$\Spec A/I$ 是 $\Spec A$ 中的一个\trRegularEmbedding。提示：证明存在 $I$ 的元素 $f_1,\dots,f_r$，构成 $k$-\trVectorSpace{} $I/(I\cap\mathfrak{m}^2)$ 的一组基。证明：$A$ 分别对 $(f_1,\dots,f_r)$ 与对 $I$ 取商，得到的都是维数为 $d$ 的\trRegular\trLocalRing。由\mbox{练习~12.1.D}，同维数的\trIntegralDomain{}之间的满同态必是同构。
\item 设 $\pi\colon X\to Y$ 是正则概形之间的一个\trClosedEmbedding。证明：$\pi$ 是一个\trRegularEmbedding。
\end{enumerate}

\mbox{练习~13.2.M}有下面这个惊人的几何推论。

\textbf{13.2.N.（练习：推广\mbox{练习~12.3.F}）}
\begin{enumerate}[label=(\alph*)]
\item 设 $W$ 是\trAlgClosedField{} $k$ 上\trPureDimension{}为 $d$ 的一个光滑\trVariety，而 $X$ 与 $Y$ 是 $W$ 的\trCodimension{}分别为 $m$ 与 $n$ 的\trPureDimensional{}子簇（可以有\trSingularPoint）。证明：$X\cap Y$ 的每个分量在 $W$ 中的\trCodimension{}至多为 $m+n$，做法如下。证明对角 $W\cong\Delta\subset W\times_k W$ 是一个\trCodimension{}为 $d$ 的\trRegularEmbedding。然后照着\mbox{练习~12.3.F}的提示把其余步骤做完。（为什么对 $W$ 的假设是光滑而非正则？）
\item 从 (a) 中去掉 $k$ \trAlgClosed{}这一要求。
\end{enumerate}

\textbf{13.2.12.}\ \emph{\trhRemark{}。}\quad 下面的例子表明，\mbox{练习~13.2.N}中的光滑性假设不能（完全）去掉。令 $W=\Spec k[w,x,y,z]/(wz-xy)$ 为光滑\trQuadricSurface{}上的锥，它是一个\trIntegral{}的三维簇。令 $X$ 为曲面 $w=x=0$，$Y$ 为曲面 $y=z=0$；二者都含于 $W$。那么 $X\cap Y$ 恰为原点，于是我们有两个\trCodimension{} $1$ 的子簇相交于一个\trCodimension{} $3$ 的子簇。（$X$ 与 $Y$ 恰是同一\trRuling{}中两条直线的\trAffine{}锥，这绝非巧合；见\mbox{练习~9.3.L}。这个例子还将在\mbox{练习~22.4.N}中再度出现。）

\textbf{13.2.13.\ \trhProposition{}。}\ ——\ \emph{光滑的 $k$-\trScheme{}都是\trReduced{}的。}

这个论证蜿蜒得出奇！

\emph{证明。}\ 设 $X$ 是一个光滑的 $k$-\trScheme，并照常令 $X_{\overline{k}}:=X\times_{\Spec k}\Spec\overline{k}$。那么 $X_{\overline{k}}$ 是光滑的（\mbox{练习~13.2.H}），故 $X_{\overline{k}}$ 在其诸\trClosedPoint{}处正则（\mbox{练习~13.2.J}），故 $X_{\overline{k}}$ 在\trClosedPoint{}处的诸\trLocalRing{}都是\trIntegralDomain（定理~13.2.11），从而既约，于是 $X_{\overline{k}}$ 在其诸\trClosedPoint{}处既约、进而是既约的（\mbox{练习~5.2.E}或~6.6.N），故 $X$ 是既约的（你看出其中的缘由吗？）。\hfill$\square$

\textbf{13.2.O.（练习）}\quad 设 $p$ 是一个维数为 $n$ 的光滑 $k$-\trScheme{} $X$ 的一个\trClosedPoint。证明：$\OO_{X,p}$ 是一个\trIntegralDomain（其维数必为 $n$）。

\textbf{13.2.14.}\ \emph{定理~13.2.11 之证明（沿 [Liu, Prop.~2.11]）。}\quad 我们对 $n$ 作归纳。$n=0$ 的情形是清楚的（维数 $0$ 的\trRegular\trLocalRing{}是域，\mbox{练习~13.2.A}）。故设 $n>0$，且结论对更小的维数成立。

设 $\mathfrak{p}$ 是素理想，且 $\dim A/\mathfrak{p}=n$。我们要证 $\mathfrak{p}=(0)$。此时 $A/\mathfrak{p}$ 是一个维数为 $n$ 的\trNoetherian\trLocalRing，它的\trZariski{}\trCotangentSpace{}是 $A$ 者的一个商空间，由 $\mathfrak{p}$ 的「方程」切出（\mbox{练习~13.1.B}）。但 $\dim\mathfrak{m}/\mathfrak{m}^2=n$，故 $\mathfrak{p}\subset\mathfrak{m}^2$。于是 $A/\mathfrak{p}$ 是一个维数为 $n$ 的\trRegular\trLocalRing。

任取 $f\in\mathfrak{m}\setminus\mathfrak{m}^2$（于是 $f\notin\mathfrak{p}$）。由\mbox{练习~12.3.I}（同时应用于 $A/(f)$ 与 $A/(\mathfrak{p},f)$），$\dim A/(f)\ge n-1$，且 $\dim A/(\mathfrak{p},f)\ge n-1$，但这两个环的\trZariski{}\trTangentSpace{}的维数都是 $n-1$，故二者都是维数为 $n-1$ 的\trRegular\trLocalRing，再由归纳假设，二者都是\trIntegralDomain。一个\trIntegralDomain{}能够作为另一个同维数\trIntegralDomain{}的商，唯一的情形便是这个商是同构（\mbox{练习~12.1.D}）。于是 $(\mathfrak{p},f)=(f)$，从而 $\mathfrak{p}\subset(f)$，从而 $\mathfrak{p}\subset f\mathfrak{p}$，从而 $\mathfrak{p}=f\mathfrak{p}$。再由\trNakayama{}引理~8.2.9，$\mathfrak{p}=(0)$，正如所愿。\hfill$\square$

\textbf{13.2.15.}\ \extrahard\ \textbf{经基变换到 $\overline{k}$，在\trClosedPoint{}处检验 $k$-\trScheme{}的\trRegularity。}

（我们将在 21.2.Y 中以另一条思路重访这些想法，所以除非你格外好奇，否则不必读本小节。）\trJacobianCriterion{}是检验\trFiniteType{} $k$-\trScheme{}在 $k$-\trValuedPoint{}处\trRegularity{}的一条极好判别法。下面的结果把它的适用面推广到更一般的\trClosedPoint。

设 $X$ 是一个\trPureDimension{}为 $n$ 的\trFiniteType{} $k$-\trScheme，而 $p\in X$ 是一个\trResidueField{}为 $k'$ 的\trClosedPoint。由\trNullstellensatz~3.2.6，$k'/k$ 是域的一个\trFiniteFieldExtension。设它是可分的。定义 $\pi\colon X_{\overline{k}}:=X\times_k\overline{k}\to X$ 为由 $\Spec\overline{k}\to\Spec k$ 经基变换而来。

\textbf{13.2.P.（练习）}
\begin{enumerate}[label=(\alph*)]
\item 设 $f(x)\in k[x]$ 是一个\trSeparable{}多项式（即 $f$ 在 $\overline{k}$ 中有互不相同的根），且不可约，于是 $k'':=k[x]/(f(x))$ 是 $k$ 的一个\trFieldExtension。证明：作为环，$k''\otimes_k\overline{k}$ 就是 $\overline{k}\times\cdots\times\overline{k}$，其中因子的个数为 $\deg f=\deg k''/k$。
\item 证明：$\pi^{-1}(p)$ 由 $\deg(k'/k)$ 个\emph{既约}点组成。
\end{enumerate}

\textbf{13.2.Q.（练习）}\quad 设 $p$ 是 $X$ 的一个\trClosedPoint，其\trResidueField{} $k'$ 是 $k$ 的 $d$ 次\trSeparable{}扩张。证明：$X_{\overline{k}}$ 在 $p$ 的全部原像 $p_1,\dots,p_d$ 处都正则，当且仅当 $X$ 在 $p$ 处正则，做法如下。
\begin{enumerate}[label=(\alph*)]
\item 化归到 $X=\Spec A$ 的情形。
\item 设 $\mathfrak{m}\subset A$ 是对应于 $p$ 的\trMaximalIdeal。将正合列
\[
\begin{tikzcd}
0 \ar[r] & \mathfrak{m} \ar[r] & A \ar[r] & k' \ar[r] & 0
\end{tikzcd}
\]
与 $\overline{k}$ 作张量（域扩张保持向量空间序列的正合性），我们便得到正合列
\[
\begin{tikzcd}
0 \ar[r] & \mathfrak{m}\otimes_k\overline{k} \ar[r] & A\otimes_k\overline{k} \ar[r] & k'\otimes_k\overline{k} \ar[r] & 0.
\end{tikzcd}
\]
证明：$\mathfrak{m}\otimes_k\overline{k}\subset A\otimes_k\overline{k}$ 是对应于 $p$ 在 $\Spec A\otimes_k\overline{k}$ 中的\trPullback{}的那个理想。验证 $(\mathfrak{m}\otimes_k\overline{k})^2=\mathfrak{m}^2\otimes_k\overline{k}$。
\item 将 $k$-\trVectorSpace{}的短正合列
\[
\begin{tikzcd}
0 \ar[r] & \mathfrak{m}^2 \ar[r] & \mathfrak{m} \ar[r] & \mathfrak{m}/\mathfrak{m}^2 \ar[r] & 0
\end{tikzcd}
\]
与 $\overline{k}$ 作张量，证明
\[
\sum_{i=1}^{d}\dim_{\overline{k}}T_{X_{\overline{k}},p_i}=d\dim_{k'}T_{X,p}.
\]
\item 利用\mbox{练习~12.1.H(a)}以及不等式 $\dim_{\overline{k}}T_{X_{\overline{k}},p_i}\ge\dim X_{\overline{k}}$ 与 $\dim_{k'}T_{X,p}\ge\dim X$（定理~12.3.12）来得出结论。
\end{enumerate}

\textbf{13.2.16.}\ \emph{\trhRemark{}。}\quad 事实上，只要有单个 $p_i$ 处的正则性，便足以推出 $p$ 处的正则性。你可以接着\mbox{练习~13.2.Q}证明这一点：先处理 $k'/k$ 是\trGalois{}扩张的情形，并得到 $\operatorname{Gal}(k'/k)$ 在 $\{p_1,\dots,p_d\}$ 上的某个传递群作用。另一条路线见\mbox{练习~21.2.Y}。

\textbf{13.2.R.（练习）}\quad 设 $k$ 是完全域，而 $X$ 是一个\trFiniteType{} $k$-\trScheme。证明：$X$ 是光滑的，当且仅当 $X$ 在所有\trClosedPoint{}处正则。特别地，定理~13.2.7(a) 成立。提示：令 $X_{\overline{k}}:=X\times_k\overline{k}$。解释为什么：$X$ 光滑当且仅当 $X_{\overline{k}}$ 光滑，当且仅当 $X_{\overline{k}}$ 在所有\trClosedPoint{}处正则，当且仅当 $X$ 在所有\trClosedPoint{}处正则。最后一步见\mbox{练习~13.2.Q}。
